Časový vývoj výdajů na aktivní politiku zaměstnanosti (\ac{APZ}) jako podílu \ac{HDP} v~ČR a~vybraných zemích \ac{EU} dokumentuje strukturální podfinancování: ČR se dlouhodobě pohybuje v~pásmu \SIrange{0.2}{0.4}{\percent} \ac{HDP}, zatímco Dánsko --- symbol flexicurity modelu --- investuje \SIrange{2.0}{2.5}{\percent} \ac{HDP} a~Německo \SIrange{0.7}{1.0}{\percent} \ac{HDP}.
Tento manko v~\ac{APZ} výdajích má přímý dopad na funkčnost trhu práce: nedostatečné rekvalifikace, slabé podpůrné systémy pro krátkodobě nezaměstnané a~omezená kapacita veřejných služeb zaměstnanosti snižují schopnost pracovníků navigovat pracovní trh po propuštění.
Flexicurity model --- na který se ČR a~vláda odvolává ve strategických dokumentech --- předpokládá trojici: flexibilní pracovní právo + štědré dávky v~nezaměstnanosti + masivní \ac{APZ}; ČR má v~zásadě jen první prvek, ostatní dva jsou výrazně podfinancovány.
Bez výrazného navýšení \ac{APZ} výdajů na alespoň \SI{1}{\percent} \ac{HDP} (německá úroveň) nemůže být flexibilizace trhu práce v~ČR legitimní: pracovníci nesou riziko flexibility, ale stát neposkytuje dostatečnou kompenzaci v~podobě přechodu a~rekvalifikace.

Časový vývoj výdajů na aktivní politiku zaměstnanosti (\ac{APZ}) jako podílu \ac{HDP} v~ČR a~vybraných zemích \ac{EU} dokumentuje strukturální podfinancování: ČR se dlouhodobě pohybuje v~pásmu \SIrange{0.2}{0.4}{\percent} \ac{HDP}, zatímco Dánsko --- symbol flexicurity modelu --- investuje \SIrange{2.0}{2.5}{\percent} \ac{HDP} a~Německo \SIrange{0.7}{1.0}{\percent} \ac{HDP}.
Tento manko v~\ac{APZ} výdajích má přímý dopad na funkčnost trhu práce: nedostatečné rekvalifikace, slabé podpůrné systémy pro krátkodobě nezaměstnané a~omezená kapacita veřejných služeb zaměstnanosti snižují schopnost pracovníků navigovat pracovní trh po propuštění.
Flexicurity model --- na který se ČR a~vláda odvolává ve strategických dokumentech --- předpokládá trojici: flexibilní pracovní právo + štědré dávky v~nezaměstnanosti + masivní \ac{APZ}; ČR má v~zásadě jen první prvek, ostatní dva jsou výrazně podfinancovány.
Bez výrazného navýšení \ac{APZ} výdajů na alespoň \SI{1}{\percent} \ac{HDP} (německá úroveň) nemůže být flexibilizace trhu práce v~ČR legitimní: pracovníci nesou riziko flexibility, ale stát neposkytuje dostatečnou kompenzaci v~podobě přechodu a~rekvalifikace.

Časový vývoj výdajů na aktivní politiku zaměstnanosti (\ac{APZ}) jako podílu \ac{HDP} v~ČR a~vybraných zemích \ac{EU} dokumentuje strukturální podfinancování: ČR se dlouhodobě pohybuje v~pásmu \SIrange{0.2}{0.4}{\percent} \ac{HDP}, zatímco Dánsko --- symbol flexicurity modelu --- investuje \SIrange{2.0}{2.5}{\percent} \ac{HDP} a~Německo \SIrange{0.7}{1.0}{\percent} \ac{HDP}.
Tento manko v~\ac{APZ} výdajích má přímý dopad na funkčnost trhu práce: nedostatečné rekvalifikace, slabé podpůrné systémy pro krátkodobě nezaměstnané a~omezená kapacita veřejných služeb zaměstnanosti snižují schopnost pracovníků navigovat pracovní trh po propuštění.
Flexicurity model --- na který se ČR a~vláda odvolává ve strategických dokumentech --- předpokládá trojici: flexibilní pracovní právo + štědré dávky v~nezaměstnanosti + masivní \ac{APZ}; ČR má v~zásadě jen první prvek, ostatní dva jsou výrazně podfinancovány.
Bez výrazného navýšení \ac{APZ} výdajů na alespoň \SI{1}{\percent} \ac{HDP} (německá úroveň) nemůže být flexibilizace trhu práce v~ČR legitimní: pracovníci nesou riziko flexibility, ale stát neposkytuje dostatečnou kompenzaci v~podobě přechodu a~rekvalifikace.

Časový vývoj výdajů na aktivní politiku zaměstnanosti (\ac{APZ}) jako podílu \ac{HDP} v~ČR a~vybraných zemích \ac{EU} dokumentuje strukturální podfinancování: ČR se dlouhodobě pohybuje v~pásmu \SIrange{0.2}{0.4}{\percent} \ac{HDP}, zatímco Dánsko --- symbol flexicurity modelu --- investuje \SIrange{2.0}{2.5}{\percent} \ac{HDP} a~Německo \SIrange{0.7}{1.0}{\percent} \ac{HDP}.
Tento manko v~\ac{APZ} výdajích má přímý dopad na funkčnost trhu práce: nedostatečné rekvalifikace, slabé podpůrné systémy pro krátkodobě nezaměstnané a~omezená kapacita veřejných služeb zaměstnanosti snižují schopnost pracovníků navigovat pracovní trh po propuštění.
Flexicurity model --- na který se ČR a~vláda odvolává ve strategických dokumentech --- předpokládá trojici: flexibilní pracovní právo + štědré dávky v~nezaměstnanosti + masivní \ac{APZ}; ČR má v~zásadě jen první prvek, ostatní dva jsou výrazně podfinancovány.
Bez výrazného navýšení \ac{APZ} výdajů na alespoň \SI{1}{\percent} \ac{HDP} (německá úroveň) nemůže být flexibilizace trhu práce v~ČR legitimní: pracovníci nesou riziko flexibility, ale stát neposkytuje dostatečnou kompenzaci v~podobě přechodu a~rekvalifikace.

\input{texparts/python/lmp_expenditure}



